\documentclass[10pt,a4paper]{amsart}
\usepackage[margin=19mm]{geometry}
\usepackage{amsmath,amssymb,mathtools}
\usepackage[scr=rsfs,cal=euler]{mathalfa}
\usepackage{xfrac}
\usepackage[unicode,hidelinks,bookmarks=false]{hyperref}
\newcommand{\ie}{i.e.\ }
\newcommand{\cf}{cf.\ }
\newcommand{\resp}{respectively\ }
\newcommand{\Subsection}{\subsection}
\providecommand{\nobreakdash}{\nobreak\hbox{-}}
\setlength{\emergencystretch}{2em}
\multlinegap0pt
\usepackage{bm}
\usepackage{bbm}
\usepackage{dsfont}
\usepackage{xspace}
\usepackage{cancel}
\usepackage{mathtools}
\usepackage{amsthm}
\makeatletter
\@ifundefined{lem}{\newtheorem{lem}{Lemma}}{}
\@ifundefined{prop}{\newtheorem{prop}{Proposition}}{}
\makeatother
\makeatletter
\expandafter\ifx\csname dedication\endcsname\relax
\newenvironment{dedication}
{\vspace{0ex}\begin{quotation}\begin{center}\begin{em}}
        {\par\end{em}\end{center}\end{quotation}}
\fi
\makeatother
\title[On the complete solutions of a generalized Lebesgue-Ramanuja]{On the complete solutions of a generalized Lebesgue-Ramanujan-Nagell equation}
\author{Kalyan Chakraborty, Azizul Hoque}
\date{Source: arXiv:2105.02832v2 (CC BY 4.0)}
\begin{document}
\maketitle

\noindent\emph{Source and licence.} On the complete solutions of a generalized Lebesgue-Ramanujan-Nagell equation. Version arXiv:2105.02832v2 (CC BY 4.0), distributed under the Creative Commons Attribution 4.0 International licence, as recorded in the arXiv licence metadata for this exact version and shown on the abstract page https://arxiv.org/abs/2105.02832v2. Full rights notice: licenses/arxiv-2105.02832-CC-BY-4.0.txt.\par\smallskip

\noindent\textit{Editorial scope.} Counted unit: the Prop whose source label is \texttt{propp} in the article (source0.tex lines 315--317, source label \texttt{propp}), delivered together with the article's own complete original proof of it (source0.tex lines 319--462, one \texttt{proof} environment of 144 lines). No mathematics is written, repaired or re-derived: statement and proof are the article's own text line for line. Required context retained uncounted: lemYU (lines 97--105); lemBH (lines 116--122); lemLE (lines 124--127); eqp1 (lines 311--313). Every definition, display and label of those lines is retained unchanged. Omitted from this package and not redistributed: 1--96, 106--115, 123--123, 128--310, 314--314, 318--318, 463--575. Nothing inside a retained excerpt is omitted or altered: every retained body is delivered exactly as the article's lines are written, so no omission receipt of any kind accompanies this package. Citations: the retained excerpts carry no citation command at all, so no bibliography entry of the article is transcribed and no reference is redistributed. Reference adaptation: none was needed and none was made; every reference inside the delivered unit resolves inside it, so no unresolved reference or citation is produced. Printed slips kept literal: no printed word, symbol, hypothesis or number of the counted statement or of its proof was altered. Layout and class: the article's own class is \texttt{amsart} with packages including amsmath, amsfonts, amssymb, hyperref, graphicx, longtable; the wrapper is the lane's pinned \texttt{amsart} editorial wrapper at 10pt on A4, which restates the theorem-like environments the retained text needs and adds the article's own macro definitions that the retained text needs (none), with extra wrapper packages (bm, bbm, dsfont, xspace, cancel, mathtools). The delivered numbering is the unit's own. Assets: no retained span contains a figure environment, a graphics-inclusion command, a PDF-page inclusion, a table or a TikZ picture; no image file is copied into this package, the package declares no asset dependency and no image file is redistributed with it. Licence: version arXiv:2105.02832v2 of this article is distributed under the Creative Commons Attribution 4.0 International licence, as recorded in the arXiv licence metadata for this exact version and shown on the abstract page https://arxiv.org/abs/2105.02832v2. A full rights notice is delivered at licenses/arxiv-2105.02832-CC-BY-4.0.txt. Characters: every retained excerpt is pure ASCII, so the delivered engine, XeLaTeX with its default Latin Modern text font, reports no missing character.
\begin{lem}\label{lemYU}
Let $d$ be a square-free positive integer and $h(-d)$ denotes the class number of $\mathbb{Q}(\sqrt{-d})$. For any odd integer $n\geq 3$ coprime to $h(-d)$, all integer solutions $(X,Y,Z)$ of the equation 
\begin{equation*}\label{eqYU}
X^2+dY^2=\lambda Z^n,~~ X,Y\geq 1, \gcd(X, dY)=1, \lambda=1,2,4,
\end{equation*} 
can be expressed as 
$$\frac{X+Y\sqrt{-d}}{\sqrt{\lambda}}=\varepsilon_1\left(\frac{a+\varepsilon_2 b\sqrt{-d}}{\sqrt{\lambda}}\right)^n,$$
where $\varepsilon_1,\varepsilon_2\in\{-1, 1\}$, and $a$ and $b$ are positive integers satisfying $\lambda Z=a^2+b^2d$ and $\gcd(a, bd)=1$.
\end{lem}

\begin{lem}\label{lemBH}
For any prime $p>13$, the Lehmer numbers $\mathcal{L}_p (\alpha, \beta) $ have primitive divisors. Further, if $p$ is a prime such that $7\leq p\leq 13$ and the Lehmer numbers $\mathcal{L}_p(\alpha, \beta)$ have no primitive divisors, then up to equivalence, the parameters $(A, B)$ of the corresponding Lehmer pair $(\alpha, \beta)$ are given by  
\begin{itemize}
\item[(i)] for $q=7, (A, B)=(1,-7), (1, -19), (3, -5), (5, -7), (13, -3), (14, -22)$;  
\item[(ii)] for $q=13, (A, B)=(1,-7)$.
\end{itemize}
\end{lem}

\begin{lem}\label{lemLE}
If $p$ is a primitive prime divisor of a Lehmer number $L_n(\alpha, \beta)$, then $p\equiv \pm 1 \pmod{n}$. Further, if $D=(\alpha^2-\beta^2)^2$ then for any odd prime $n$, $p\equiv \left(\frac{D}{p}\right)\pmod n$, where $\left(\frac{D}{p}\right)$
denotes the Legendre symbol.
\end{lem}

\begin{equation}\label{eqp1}
x^2+17^k41^\ell 59^m =\lambda y^p,~~x\geq 1, y>1, \gcd(x,y)=1, k, \ell, m\geq 0, p\geq 5,
\end{equation}

\begin{prop}\label{propp}The only integer solution of \eqref{eqp1} is $(x,y,\lambda,k,\ell,m,p)=(38,5,1,0,2,0,5)$.

\end{prop}

\begin{proof}
Suppose that $k=2k_1+k_2, \ell=2\ell_1+\ell_2$ and $m=2m_1+m_2$, where $k_1,\ell_1, m_1\geq 0$ are integers and $k_2,\ell_2,m_2\in\{0,1\}$. Then 
$17^k41^\ell 59^m=dz^2$, where $d=17^{k_2}41^{\ell_2} 59^{m_2}$ and $z=17^{k_1}41^{\ell_1} 59^{m_1}$. 
Therefore \eqref{eqp1} becomes
$$x^2+dz^2=\lambda y^p, ~~x\geq 1, y>1, \gcd(x,dz)=1, k, \ell, m\geq 0, p\geq 5, \lambda =1, 2,4.$$
Using MAGMA, we see that $h(-d)\in\{1,3,4,8,72\}$, where $h(-d)$ denotes the class number of $\mathbb{Q}(\sqrt{-d})$.
Thus by Lemma \ref{lemYU}, we have
\begin{equation}\label{eqp2}
\frac{x+z\sqrt{-d}}{\sqrt{\lambda}}=\varepsilon_1\left(\frac{a+\varepsilon_2 b\sqrt{-d}}{\sqrt{\lambda}}\right)^p,
\end{equation} 
where $\varepsilon_1, \varepsilon_2\in\{-1, 1\}$, and $a$ and $b$ are positive integers satisfying
\begin{equation}\label{eqp3}
a^2+b^2d=\lambda y.
\end{equation}
Clearly, $b$ is odd. Since $y$ is odd, so that $a$ is even ( resp. odd) when $\lambda =1$ (resp. $\lambda = 2, 4$). 

Assume that $$\alpha:=\frac{a+\varepsilon_2 b\sqrt{-d}}{\sqrt{\lambda}},~~\beta:=\frac{a-\varepsilon_2 b\sqrt{-d}}{\sqrt{\lambda}}.$$
Clearly, both $\alpha$ and $\beta$ are algebraic integers satisfying $\gcd((\alpha+\beta)^2, \alpha\beta)=1$. Also, $\alpha/\beta$ is a root of $\lambda yZ^2-2(a^2-b^2d)+\lambda y=0$, and hence $\alpha/\beta$ is not a root of unity. Therefore, $(\alpha, \beta)$ is  Lehmer pair with parameters $(4a^2/\lambda, -4b^2d/\lambda)$. 

If $\mathcal{L}_n$ is the $n$-th Lehmer number for the pair $(\alpha, \beta)$, then 
\begin{equation}\label{eqp4}
|\mathcal{L}_p(\alpha, \beta)|=\frac{z}{b}=\frac{17^{k_1}41^{\ell_1} 59^{m_1}}{b}.
\end{equation}
Let $q$ be a primitive prime divisor of $\mathcal{L}(\alpha, \beta)$. Then by \eqref{eqp4}, $q\in\{17,41,59\}$. Utilizing Lemma \ref{lemLE}, we see that $q\ne 17$ since $p\geq 5$. Also, 
$$q=\begin{cases}
41 & \text{ when } p=5,7;\\
59 & \text{ when } p=5, 29.
\end{cases}
$$
We first consider $q=59$. Then by the definition of primitive divisor, $q\nmid (\alpha^2-\beta^2)^2=-16a^2b^2d/\lambda^2$. This gives that  $(k_2, \ell_2, m_2)\in\{(0,0,0), (0,1,0), (1,0,0), (1,1,0)\}$, which implies $d\in\{1,41,17, 697\}$. Again since $\left(\frac{-4b^2d}{59}\right)=\left(\frac{-d}{59}\right)=-1$, so that $59\equiv 1\pmod p$ implies $p=5$.

If $q=41$, then as before we get  $(k_2, \ell_2, m_2)\in\{(0,0,0), (0,0,1), (1,0,0), (1,0,1)\}$ and thus $d\in\{1,59, 17, 1003\}$. Again since $\left(\frac{-4b^2d}{41}\right)=\left(\frac{-d}{41}\right)=1$ or $-1$ according as $d=1,59$ or $d=17, 1003$, so that $41\equiv \pm 1\pmod p$ gives $(p, d)\in\{(7,17),(7,1003), (5,1), (5,59)\}$. We now divide the proof into several cases depending on the values of $p$ and $q$.
%%%%%%%%%%%%%%
\subsection*{Case I:  $p>7$} 
From the above, it follows that $\mathcal{L}_p(\alpha, \beta)$ has no primitive divisors for $p>7$. This contradicts to the primitive divisor theorem for Lehmer sequences, which states that, if $p\geq 3$, then $\mathcal{L}_p(\alpha, \beta)$ has a primitive prime divisor with a few exceptional pairs $(\alpha, \beta)$. For any prime $p> 7$, these exceptional pairs are given by Lemma \ref{lemBH} in terms of their parameters. Therefore using Lemma \ref{lemBH} we get $(p,4a^2/\lambda, 4b^2d/\lambda)=(13,1,7)$, which further implies that $d=7$. This is not possible as the only prime divisors of $d$ are $17, 41, 59$. Therefore, \eqref{eqp1} has no integer solutions for $p>7$. 
%%%%%%%%%%%%
\subsection*{Case II:  $p=7$} 
In this case, the only primitive divisor of $\mathcal{L}_7(\alpha, \beta)$ is $41$, and accordingly  $d=17, 1003$. Equating the imaginary parts in \eqref{eqp2}, we get 
\begin{equation}\label{eqp5}
\lambda^317^{k_1}41^{\ell_1}59^{m_1}=\varepsilon b(7a^6-35a^4b^2d+21a^2b^4d^2-b^6d^3),
\end{equation}
where $\varepsilon=\varepsilon_1\varepsilon_2=\pm 1$ and $d=17, 1003$. Also by the definition of primitive divisor, $41\nmid (\alpha^2-\beta^2)^2=-16a^2b^2d/\lambda$, we have $b\ne 41$. Since $b$ is odd and $\gcd(a, b)=1$, so that \eqref{eqp5} gives 
$$b=1,17^{k_1},59^{m_1}, 17^{k_1}59^{m_1}.$$

We first consider $b=1$. Then \eqref{eqp5} reduces to
$$\lambda^317^{k_1}41^{\ell_1}59^{m_1}=\varepsilon (7a^6-35a^4d+21a^2d^2-d^3).$$
 Assume that $k_1=2k_4+k_3, \ell_1=2\ell_4+\ell_3, m_1=2m_4+m_3$ with $k_3, \ell_3, m_3\in\{0, 1\}$ and then set $D:=\varepsilon\lambda 17^{k_3}41^{\ell_3}59^{m_3}$ to get
$$D(\lambda 17^{k_4}41^{\ell_4}59^{m_4})^2=7a^6-35a^4d+21a^2d^2-d^3.$$
Multiplying both sides by $7^2D^3$, and then putting $(X, Y):=(7Da^2, 7D^2\lambda 17^{k_4}41^{\ell_4}59^{m_4})$, we arrive at the following:
$$Y^2=X^3-35dDX^2+147d^2D^2X-49d^3D^3.$$
This defines $96$ elliptic curves, where we use \texttt{IntegralPoints} subroutine of MAGMA to compute all integer points. We first remove all integer points $(X,Y)$ satisfying $XY=0$ or $7\nmid X$ since $abd\ne0$ and $7\mid \gcd(X,Y)$. We then check $a^2=X/{7D}$ for the remaining points $(X,Y)$; but none of these shows that $a$ is an integer. Thus, none of these integer points lead to an integer solutions of \eqref{eqp1}. 

Assume that $b=17^{k_1}$. Then \eqref{eqp5} reduces to
$$\lambda^341^{\ell_1}59^{m_1}=\varepsilon (7a^6-35a^4b^2d+21a^2b^4d^2-b^6d^3).$$
As before, we consider $D_1:=\varepsilon\lambda 41^{\ell_3}59^{m_3}$ for $ \ell_3, m_3\in\{0, 1\}$. Then the above equation reduces to
$$D_1(\lambda 41^{\ell_4}59^{m_4})^2=7a^6-35a^4b^2d+21a^2b^4d^2-b^6d^3,$$
 where $\ell_1=2\ell_4+\ell_3, m_1=2m_4+m_3$.
We divide both sides by $b^6$ and then multiply by $7^2D_1^3$ to get
$$Y^2=X^3-35dD_1X^2+147d^2D_1^2X-49d^3D_1^3,$$
where $X=7D_1a^2/b^2$ and $Y=7D_1\lambda 41^{\ell_4}59^{m_4}/b^3$.
As in the previous case, we use \texttt{IntegralPoints} subroutine of MAGMA to compute all $\{17\}$-integral points on the above $48$ elliptic curves. We eliminate all integer points $(X,Y)$ satisfying $XY=0$ or $7\nmid X$ since $abd\ne0$ and $7\mid \gcd(X,Y)$. For the remaining points $(X,Y)$, we examine $a^2=Xb^2/{7D_1}$. However, none of these gives integer value to $a$, and hence none of these integer points lead to an integer solution of \eqref{eqp1}. 

Let $b=59^{m_1}$. Then \eqref{eqp5} becomes
$$\lambda^317^{k_1}41^{\ell_1}=\varepsilon (7a^6-35a^4b^2d+21a^2b^4d^2-b^6d^3).$$
As before, put $D_2:=\varepsilon\lambda 17^{k_3}41^{\ell_3}$ for $k_3, \ell_3, m_3\in\{0, 1\}$. Then we have
$$D_2(\lambda 17^{k_4}41^{\ell_4})^2=7a^6-35a^4b^2d+21a^2b^4d^2-b^6d^3,$$
 where $k_1=2k_4+k_3, \ell_1=2\ell_4+\ell_3$.
We divide both sides by $b^6$ and then multiply by $7^2D_1^3$ to get
$$Y^2=X^3-35dD_2X^2+147d^2D_2^2X-49d^3D_2^3,$$
where $X=7D_2a^2/b^2$ and $Y=7D_2\lambda 17^{k_4}41^{\ell_4}/b^3$.
In this case too, we use \texttt{IntegralPoints} subroutine of MAGMA to compute all $\{59\}$-integral points on the above $48$ elliptic curves.  As in the last case, we see that one  none of these integer points lead to an integer solution of \eqref{eqp1}. 


Finally consider $b=17^{k_1}59^{m_1}$. Then \eqref{eqp5} reduces to
$$\lambda^341^{\ell_1}=\varepsilon (7a^6-35a^4b^2d+21a^2b^4d^2-b^6d^3).$$
As before, putting $D_3:=\varepsilon\lambda 41^{\ell_3}$ for $\ell_3\in\{0, 1\}$ in the above equation, we get 
$$D_3(\lambda 41^{\ell_4})^2=7a^6-35a^4b^2d+21a^2b^4d^2-b^6d^3,$$
 where $\ell_1=2\ell_4+\ell_3$.
We now divide both sides by $b^6$ and then multiply by $7^2D_3^3$ to arrive at 
$$Y^2=X^3-35dD_3X^2+147d^2D_3^2X-49d^3D_3^3,$$
where $X=7D_3a^2/b^2$ and $Y=7D_3\lambda 41^{\ell_4}/b^3$.
As in previous cases, we use \texttt{IntegralPoints} subroutine of MAGMA to compute all $\{17,59\}$-integral points on the above $24$ elliptic curves. Utilizing the same technique as before, we can conclude that none of these integer points lead to an integer solution of \eqref{eqp1}. 

%%%%%%%%%%%%
\subsection*{Case III:  $p=5$} 
In this case, both $41$ and $59$ are primitive divisors of $\mathcal{L}_5(\alpha, \beta)$, and 
$(q, d)=(41, 1), (41, 59), (59, 1), (59, 17), (59, 41), (59, 1003)$. 
Equating the imaginary parts in \eqref{eqp2}, we get 
\begin{equation}\label{eqp6}
\lambda^217^{k_1}41^{\ell_1}59^{m_1}=\varepsilon b(5a^4-10a^2b^2d+b^4d^2),
\end{equation}
where $\varepsilon=\varepsilon_1\varepsilon_2=\pm 1$ . In case of $q=41$,  by the definition of primitive divisor, $b\ne 41$. Thus \eqref{eqp6} gives  
$$b=1,17^{k_1},59^{m_1}, 17^{k_1}59^{m_1}.$$
Similarly, when $q=59$,  we get (from \eqref{eqp6})
$$b=1,17^{k_1}, 41^{\ell_1}, 17^{k_1}41^{\ell_1}.$$

If $b=1$, then $d=1,17,41,59,1003$ and \eqref{eqp6} reduces to
$$\lambda^217^{k_1}41^{\ell_1}59^{m_1}=\varepsilon (5a^4-10a^2d+d^2).$$
Let $k_1=2k_4+k_3, \ell_1=2\ell_4+\ell_3, m_1=2m_4+m_3$ with $k_3,\ell_3, m_3\in\{0,1\}$ and set $D:=\varepsilon 17^{k_3}41^{\ell_3}59^{m_3}$. Then the above equation becomes 
$$D(\lambda 17^{k_4}41^{\ell_4}59^{m_4})^2=5a^4-10a^2d+d^2.$$
We now multiply both sides by $D$, and then put $(X, Y):=(a, D\lambda 17^{k_4}41^{\ell_4}59^{m_4})$ to get
$$Y^2=5DX^4-10DdX^2+Dd^2.$$
The above equation gives $80$ quartic curves each for each value of $Dd$. Here, we use \texttt{IntegralQuarticPoints} subroutine of MAGMA to compute all integral points on these curves. As in Case II, we first eliminate the integer points $(X,Y)$ such that $XY=0$ since $abd\ne0$. For the remaining points, we have $(X,Y, D)=(1,-2,-1), (2,41,41), (2,11,1),$ $ (12,149,1)$. Clearly, $(X,Y, D)=(1,-2,-1)$ gives $(x,y,\lambda,k,\ell,m)=(1,1,2,0,0,0)$. On the other hand, $(X,Y,D)=(2,41,41)$ gives $(a,\lambda, k,\ell, m)=(2,1,0,2,0)$. Here, $d=1$ and hence by \eqref{eqp3} $y=5$, which further implies $x=38$. Therefore $(x,y,\lambda, k, \ell, m, p)=(38,5,1,0,2,0,5)$ is a solution of \ref{eqp1}. Also, since $2,17,41,59$ are the only possible prime factors of $Y$, so that $(X,Y,D)=(2,11,1), (12,149,1)$ are not possible. 



We now consider the case $b=17^{k_1}$.  Then \eqref{eqp6} becomes 
$$\lambda^241^{\ell_1}59^{m_1}=\varepsilon (5a^4-10a^2b^2d+b^4d^2).$$
Here, $d=1,17,41,59,1003$.
Let $D_1:=\varepsilon 41^{\ell_3}59^{m_3}$ for $k_3, \ell_3, m_3\in\{0, 1\}$, and divide both sides by $b^4$, we get
$$D_1(\lambda 41^{\ell_4}59^{m_4}/b^2)^2=5(a/b)^4-10(a/b)^2d+d^2,$$
 where $ \ell_1=2\ell_4+\ell_3, m_1=2m_4+m_3$.
We set $(X, Y):=(a/b, \lambda 41^{\ell_4}59^{m_4}/b^2)$ to arrive at the following:
$$D_1Y^2=5X^4-10dX^2+d^2.$$
We use \texttt{SIntegralLjunggrenPoints} subroutine of MAGMA to compute all $\{17\}$-integral points on the $40$ quartic curves defined by the above equation. As in the previous cases, we get $(X,Y, D)=(1,2,-1), (2,1,41), (2,11,1), (12,149,1)$. 
As before, these give $(x,y,\lambda,k,\ell,m,p)=(1,1,2,0,0,0,5), (38,5,1,0,2,0,5)$.

Assume that $b=41^{\ell_1}$.  Then $d=1, 59$ and \eqref{eqp6} reduces to
$$\lambda^217^{k_1}59^{m_1}=\varepsilon (5a^4-10a^2b^2d+b^4d^2).$$
Let $D_2:=\varepsilon 17^{k_3}59^{m_3}$ for $k_3, m_3\in\{0, 1\}$, and divide both sides by $b^4$, we get
$$D_2(\lambda 17^{k_4}59^{m_4}/b^2)^2=5(a/b)^4-10(a/b)^2d+d^2,$$
where $ k_1=2k_4+k_3, m_1=2m_4+m_3$.
Putting $(X, Y):=(a/b, \lambda 17^{k_4}59^{m_4}/b^2)$, we get
$$D_2Y^2=5X^4-10dX^2+d^2.$$
Here too we use \texttt{SIntegralLjunggrenPoints} subroutine of MAGMA to compute all $\{41\}$-integral points on the $16$ quartic curves defined by the last equation. As before, we get $(X,Y, D)=(1,2,-1), (12,149,1)$, which lead to  $(x,y,\lambda,k,\ell,m,p)=(1,1,2,0,0,0,5)$.

For $b=59^{m_1}$,  \eqref{eqp6} becomes 
$$\lambda^217^{k_1}41^{\ell_1}=\varepsilon (5a^4-10a^2b^2d+b^4d^2).$$
As in the previous cases, by setting $D_3:=\varepsilon 17^{k_3}41^{\ell_3}$, this equation can written as
$$D_3(\lambda 17^{k_4}41^{\ell_4}/b^2)^2=5(a/b)^4-10(a/b)^2d+d^2,$$
where $ k_1=2k_4+k_3, \ell_1=2\ell_4+\ell_3$.
Assume that $(X, Y):=(a/b, \lambda 17^{k_4}41^{\ell_4}/b^2)$. Then we get
$$D_3Y^2=5X^4-10dX^2+d^2.$$
As before, we get $(X,Y, D)=(1,2,-1), (2,1,41), (2,11,1)$. The solutions corresponding to these points are already discussed.

Finally for $b=41^{\ell_1}59^{m_1}$, \eqref{eqp6} implies to
$$\lambda^217^{k_1}=\varepsilon (5a^4-10a^2b^2d+b^4d^2).$$
We put $D_4:=\varepsilon 17^{k_3}$ for $k_3\in\{0, 1\}$, and divide both sides by $b^4$ to get
$$D_4(\lambda 17^{k_4}/b^2)^2=5(a/b)^4-10(a/b)^2d+d^2,$$
where $ k_1=2k_4+k_3$.
Writing $(X, Y):=(a/b, \lambda 17^{k_4}$, we get
$$D_3Y^2=5X^4-10dX^2+d^2.$$
Using \texttt{SIntegralLjunggrenPoints} subroutine of MAGMA, we compute all $\{41,59\}$-integral points on the $20$ quartic curves defined by the last equation. Here, we get $(X,Y, D)=(1,2,-1), (2,11,1), (12, 149,1)$, which are already treated in the previous cases. This completes the proof.
\end{proof}

\end{document}
